\documentclass{easychair}

\usepackage{amsmath}
\usepackage{amssymb}
\usepackage{amsfonts}
\usepackage{algorithm}
\usepackage{algorithmic}
\usepackage{float}
\usepackage{graphicx}
\usepackage{url}
\usepackage{booktabs}
\usepackage{multirow}

\title{Toward Quantum-Resilient eSIM Provisioning: Integrating ML-KEM into Consumer Remote SIM Provisioning}

\author{
  Jhury Kevin Lastre \and
  Hoseok Kwon \and
  Bonam Kim \and
  Ilsun You\thanks{Corresponding author}
}
\institute{
  Kookmin University, Seoul, South Korea \\
  \email{lavelliane@kookmin.ac.kr, hoseok1997@kookmin.ac.kr, kimbona@kookmin.ac.kr, isyou@kookmin.ac.kr}
}
\authorrunning{Lastre, Kwon, Kim and You}
\titlerunning{Toward Quantum-Resilient eSIM Provisioning}

\begin{document}

\maketitle

\begin{abstract}
Consumer Remote SIM Provisioning (RSP), standardized through GSM Association (GSMA) specifications SGP.22 and SGP.24, establishes secure embedded Subscriber Identity Module (eSIM) profile management using Public Key Infrastructure (PKI) based trust models between Subscription Manager Data Preparation (SM-DP+), embedded Universal Integrated Circuit Card (eUICC), and Local Profile Assistant (LPA) components. As quantum computing threatens classical cryptographic primitives such as Rivest-Shamir-Adleman (RSA) and Elliptic Curve Diffie-Hellman (ECDH) through Shor's algorithm, telecommunications infrastructure requires transition strategies that preserve interoperability with existing GSMA-trusted Certificate Authorities while achieving quantum resilience. This work presents a Post-Quantum Cryptography (PQC) integration strategy that introduces quantum-safe mechanisms at the algorithmic layer using liboqs to enable hybrid ECDH plus Module Lattice Key Encapsulation Mechanism (ML-KEM) key exchanges. The proposed migration maintains the GSMA certificate hierarchy and SGP.22 protocol structure, allowing PQC-capable eUICC implementations to process quantum-resistant key material with backward compatibility for legacy systems. Experimental evaluation on a virtual testbed demonstrates that hybrid key agreement introduces 2,272 bytes of additional key material and 0.232 milliseconds of computational overhead per provisioning session, resulting in 7.5x bandwidth increase while achieving NIST security level 3 quantum resistance.
\end{abstract}

\textbf{Keywords:} Remote SIM Provisioning, eSIM, Post-Quantum Cryptography, ML-KEM, Hybrid Key Exchange, Quantum Resistance

\section{Introduction}

Remote SIM Provisioning (RSP), formalized through the GSM Association (GSMA) specifications SGP.22 and SGP.24 \cite{GSMA_SGP22,GSMA_SGP24}, establishes the standardized framework for mobile network operators to securely provision and manage embedded Subscriber Identity Module (eSIM) profiles over-the-air. The architecture guarantees end-to-end confidentiality and cryptographic integrity across communications between the Subscription Manager Data Preparation (SM-DP+) server, the embedded Universal Integrated Circuit Card (eUICC), and the Local Profile Assistant (LPA) component through a Public Key Infrastructure (PKI) based trust model. With over 1.2 billion eSIM-capable devices deployed globally as of 2024, RSP has become a critical component of telecommunications infrastructure, enabling seamless operator switching, international roaming, and machine-to-machine connectivity.

The security of RSP relies fundamentally on classical cryptographic primitives, specifically Elliptic Curve Digital Signature Algorithm (ECDSA) with NIST P-256 for authentication and Elliptic Curve Diffie-Hellman (ECDH) for session key establishment. However, these primitives are vulnerable to quantum computational attacks. Shor's algorithm \cite{NIST_PQC2024}, executable on sufficiently large quantum computers, can solve both the discrete logarithm problem and integer factorization in polynomial time, thereby breaking RSA, ECDSA, and ECDH. Current projections suggest that cryptographically relevant quantum computers capable of breaking 2048-bit RSA could emerge within 10 to 15 years. This timeline presents an urgent threat when considering the long-term sensitivity of telecommunications credentials and the "store-now-decrypt-later" attack model, where adversaries harvest encrypted eSIM profile data today for future decryption once quantum computers become available.

Given the critical nature of subscriber credentials, operator authentication keys, and profile encryption keys transmitted during RSP provisioning, the transition to Post-Quantum Cryptography (PQC) is imperative. However, this migration presents significant challenges. The GSMA trust infrastructure involves a complex hierarchy of Certificate Authorities (CAs), with EUM (eUICC Manufacturer) and SM (Subscription Manager) certificates issued by GSMA-trusted root CAs. Any migration strategy must preserve this established PKI ecosystem to maintain interoperability with the existing deployed base of millions of eUICC devices and hundreds of SM-DP+ servers worldwide. Furthermore, the constrained nature of eUICC hardware, with limited memory (typically 256KB to 1MB) and processing power, imposes strict requirements on the size and computational cost of cryptographic operations.

\subsection{Research Questions and Contributions}

This work addresses three fundamental research questions:

\begin{enumerate}
\item \textbf{How can PQC be integrated into SGP.22 without disrupting the existing GSMA PKI hierarchy?} We demonstrate that quantum-resistant key exchange can be achieved through algorithmic-layer modifications while preserving certificate structures and trust chains.

\item \textbf{What is the performance overhead of hybrid key agreement in resource-constrained eSIM environments?} Through experimental evaluation on a virtual testbed, we quantify the bandwidth, memory, and computational costs of ML-KEM-768 integration.

\item \textbf{What deployment challenges arise in migrating real-world RSP implementations to support PQC?} We document practical issues encountered during implementation, including buffer management, TLV encoding complexities, and backward compatibility mechanisms.
\end{enumerate}

Our primary contributions are:

\begin{itemize}
\item A hybrid key exchange protocol that combines classical ECDH with ML-KEM-768 for quantum-resistant session key derivation in SGP.22, maintaining backward compatibility with legacy eUICCs.

\item ASN.1 protocol extensions for PrepareDownload and InitialiseSecureChannel messages to transport ML-KEM public keys and ciphertexts without modifying the core SGP.22 message structure.

\item An implementation and experimental evaluation demonstrating 7.5x bandwidth overhead with negligible (0.232 ms) computational cost per provisioning session.

\item Analysis of real-world deployment challenges including buffer overflow vulnerabilities, long-form TLV encoding requirements, and secure key management in constrained environments.
\end{itemize}

The remainder of this paper is organized as follows. Section 2 provides background on GSMA RSP architecture and PQC fundamentals. Section 3 describes our migration architecture and design principles. Section 4 details the implementation of hybrid key agreement on both eUICC and SM-DP+ sides. Section 5 describes our experimental testbed. Section 6 presents evaluation results with performance analysis. Section 7 discusses deployment challenges and standardization considerations. Section 8 surveys related work, and Section 9 concludes with future research directions.

\section{Background}

\subsection{GSMA Remote SIM Provisioning Architecture}

The GSMA RSP architecture, standardized in SGP.22 for consumer devices, defines secure procedures for over-the-air eSIM profile installation. The architecture comprises three primary entities: the SM-DP+ server operated by the mobile network operator or service provider, the eUICC embedded in the end-user device, and the LPA software component running on the device's operating system that facilitates communication between the eUICC and SM-DP+.

The profile download and installation flow proceeds through four main phases as shown in Table \ref{tab:protocol-phases}. During Phase 1 (Authentication), the eUICC and SM-DP+ perform mutual authentication using their respective ECDSA certificates issued by GSMA-trusted CAs. The eUICC generates a challenge and the SM-DP+ responds with its certificate chain. The eUICC verifies this chain against its trusted root CA public keys, then signs its own challenge response using its ECDSA private key. Phase 2 (PrepareDownload) establishes session key material. The eUICC generates an ephemeral ECDH key pair (otPK.EUICC.ECKA, otSK.EUICC.ECKA) and transmits the public key to the SM-DP+ along with metadata about its capabilities and available memory. Phase 3 (Bound Profile Package Download) delivers the encrypted profile. The SM-DP+ generates its own ephemeral ECDH key pair, performs key agreement to derive shared secrets, and uses these to derive session encryption keys (KEK) and message authentication keys (KM). The Bound Profile Package (BPP) contains the encrypted profile data, SIM application code, and personalization data. Phase 4 (Installation) involves the eUICC decrypting the BPP, verifying MAC tags, installing profile elements, and returning a signed ProfileInstallationResult.

\begin{table}[h]
\centering
\caption{SGP.22 Profile Download Protocol Phases}
\label{tab:protocol-phases}
\begin{tabular}{@{}lll@{}}
\toprule
\textbf{Phase} & \textbf{Protocol Messages} & \textbf{Cryptographic Operations} \\
\midrule
1. Authentication & ES9+.InitiateAuthentication & ECDSA signature verification \\
 & ES10b.AuthenticateServer & Certificate chain validation \\
 & ES9+.AuthenticateClient & ECDSA signature generation \\
\midrule
2. PrepareDownload & ES10b.PrepareDownload & ECDH keypair generation \\
 & ES9+.GetBoundProfilePackage & Ephemeral key transmission \\
\midrule
3. BPP Download & ES8+.InitialiseSecureChannel & ECDH key agreement \\
 & ES8+.ConfigureISDP & Session key derivation (KEK, KM) \\
 & ES8+.StoreMetadata & Profile decryption (AES-128-CBC) \\
 & ES8+.LoadProfileElements & MAC verification (CMAC) \\
\midrule
4. Installation & ES10b.ProfileInstallationResult & Installation verification \\
 & ES9+.HandleNotification & Result signature (ECDSA) \\
\bottomrule
\end{tabular}
\end{table}

The security of this protocol relies on the confidentiality and authenticity provided by ECDSA P-256 signatures and ECDH P-256 key agreement. The session keys (KEK and KM) are derived using a Key Derivation Function (KDF) based on SHA-256, taking the ECDH shared secret as input along with domain-specific labels. The KEK is used for AES-128-CBC encryption of profile data, while KM is used for CMAC-based message authentication to prevent tampering.

\subsection{Post-Quantum Cryptography and ML-KEM}

The National Institute of Standards and Technology (NIST) initiated a PQC standardization project in 2016 to identify quantum-resistant cryptographic algorithms \cite{NIST_PQC2024}. After multiple rounds of evaluation, NIST selected CRYSTALS-Kyber (now standardized as ML-KEM, Module Lattice Key Encapsulation Mechanism) for key establishment in 2024. ML-KEM is based on the hardness of the Module Learning With Errors (MLWE) problem over polynomial rings, which remains computationally intractable even for quantum adversaries.

ML-KEM operates as a Key Encapsulation Mechanism (KEM) rather than a traditional key agreement protocol like Diffie-Hellman. The security model differs fundamentally: in a KEM, one party (the encapsulator) generates a shared secret and encrypts it under the other party's public key, producing a ciphertext. The recipient (the decapsulator) uses their private key to decrypt this ciphertext and recover the shared secret. This asymmetric approach contrasts with ECDH, where both parties contribute equally to derive a shared secret through scalar multiplication on an elliptic curve.

ML-KEM-768 specifically provides NIST security level 3, which offers 192 bits of quantum security, comparable to the classical security of AES-192. The key sizes are significantly larger than classical equivalents: the public key is 1,184 bytes, the secret key is 2,400 bytes, the ciphertext is 1,088 bytes, and the shared secret is 32 bytes. These larger sizes pose challenges for bandwidth-constrained protocols and memory-limited devices. However, the computational performance of ML-KEM is competitive with classical cryptography. Keypair generation and encapsulation operations can be performed in under 1 millisecond on modern processors, and decapsulation is even faster.

\subsection{Threat Model and Quantum Adversary}

Our threat model considers an adversary with access to a cryptographically relevant quantum computer capable of executing Shor's algorithm to break ECDH and ECDSA. We assume the adversary can intercept and store all RSP protocol messages transmitted between the eUICC and SM-DP+ over potentially untrusted networks (e.g., public Wi-Fi, cellular data). The adversary's goal is to decrypt the BPP and extract sensitive data including the subscriber's International Mobile Subscriber Identity (IMSI), Ki authentication key, operator credentials, and personalization data.

Under the "store-now-decrypt-later" attack scenario, the adversary captures encrypted profile data transmitted today and stores it for future cryptanalysis. Once a sufficiently powerful quantum computer becomes available, the adversary can:

\begin{enumerate}
\item Break the ECDSA signatures on authentication messages to forge certificates or impersonate entities.
\item Solve the ECDH discrete logarithm problem to recover the ephemeral private key otSK.EUICC.ECKA.
\item Compute the shared secret and derive the session keys (KEK, KM).
\item Decrypt the BPP and extract all sensitive profile data.
\end{enumerate}

This threat is particularly acute for eSIM profiles because subscriber credentials have long-term value. A compromised Ki key enables the adversary to authenticate as the legitimate subscriber indefinitely, intercept communications, and commit fraud. Furthermore, the deployment lifetime of eUICC devices (often 10+ years for IoT applications) means that devices provisioned today using only classical cryptography will remain vulnerable throughout their operational life.

Our mitigation strategy focuses on protecting the session key establishment phase (Phase 2 and 3) by replacing pure ECDH with a hybrid ECDH plus ML-KEM key agreement. This ensures that even if a quantum adversary breaks the ECDH component in the future, the ML-KEM component remains secure, preserving the confidentiality of the session keys and consequently the BPP. We defer the migration of signature algorithms (ECDSA to ML-DSA) to future work, as signatures provide authentication but not long-term confidentiality. Breaking a signature allows impersonation but does not retroactively compromise past sessions.

\section{Migration Architecture}

\subsection{Design Principles}

Our PQC migration strategy is guided by four key principles that prioritize practical deployment and ecosystem compatibility:

\textbf{Principle 1: Preserve GSMA PKI Infrastructure.} The existing GSMA certificate hierarchy must remain intact. We do not replace or modify EUM and SM certificates, nor do we require changes to GSMA root CA operations. PQC is integrated solely at the session key establishment layer, orthogonal to the PKI.

\textbf{Principle 2: Ensure Backward Compatibility.} Legacy eUICCs that do not support PQC must continue to function without modification. The protocol must gracefully fall back to classical ECDH when either the eUICC or SM-DP+ does not advertise PQC capabilities. This enables incremental deployment across the ecosystem.

\textbf{Principle 3: Adopt Hybrid Cryptography During Transition.} Rather than immediately replacing ECDH with pure ML-KEM, we employ a hybrid approach that combines both mechanisms. This provides defense-in-depth: the security level is the maximum of the two schemes, so the protocol remains secure even if one algorithm is unexpectedly broken.

\textbf{Principle 4: Minimize Protocol Disruption.} PQC integration should require minimal changes to SGP.22 message structures. We avoid redesigning the protocol or introducing new message types. Instead, we extend existing messages (PrepareDownloadResponse and InitialiseSecureChannelRequest) with optional TLV fields to carry ML-KEM key material.

\subsection{Hybrid Key Agreement Protocol}

Figure \ref{fig:architecture} illustrates our migration architecture. The protocol flow maintains the four-phase structure of SGP.22 but augments Phase 2 and Phase 3 with hybrid key agreement operations.

\begin{figure}[h]
  \centering
  \includegraphics[width=0.85\columnwidth]{output/poster-diagram.png}
  \caption{Migration architecture showing hybrid key agreement flow between SM-DP+, LPA, and eUICC. The eUICC generates both ECDH and ML-KEM keypairs during PrepareDownload. The SM-DP+ performs hybrid key agreement combining ECDH shared secret with ML-KEM encapsulation to derive quantum-resistant session keys.}
  \label{fig:architecture}
\end{figure}

During PrepareDownload, the eUICC advertises its PQC capabilities through the EUICCInfo2 structure. If the eUICC supports ML-KEM-768, it generates both the classical ECDH keypair (otPK.EUICC.ECKA, otSK.EUICC.ECKA) and an ML-KEM-768 keypair (pk\_kem.EUICC, sk\_kem.EUICC). The PrepareDownloadResponse message includes both the ECDH public key (as per standard SGP.22) and the ML-KEM public key as an optional extension. The SM-DP+ receives this response and checks for the presence of the ML-KEM public key. If found, the SM-DP+ enters hybrid mode.

During BPP generation, the SM-DP+ generates its own ECDH keypair (otPK.DP.ECKA, otSK.DP.ECKA) and performs classical ECDH key agreement to obtain shared secret Z\_ecdh. Additionally, the SM-DP+ invokes ML-KEM-768 encapsulation using the eUICC's ML-KEM public key to generate shared secret Z\_mlkem and ciphertext ct\_kem. The InitialiseSecureChannelRequest message (the first command in the BPP) includes both the SM-DP+ ECDH public key and the ML-KEM ciphertext as extensions.

When the eUICC receives the InitialiseSecureChannelRequest, it performs classical ECDH with the SM-DP+ public key to derive Z\_ecdh, then invokes ML-KEM-768 decapsulation on the ciphertext using its ML-KEM secret key to recover Z\_mlkem. With both shared secrets available, the eUICC applies a hybrid Key Derivation Function (KDF) to derive the final session keys (KEK and KM). The SM-DP+ performs the same hybrid KDF using its copies of Z\_ecdh and Z\_mlkem, resulting in both parties deriving identical session keys.

If either party does not support ML-KEM or does not advertise the capability, the protocol falls back to classical-only mode. The eUICC omits the ML-KEM public key from PrepareDownloadResponse, the SM-DP+ omits the ciphertext from InitialiseSecureChannelRequest, and both parties use the standard SGP.22 KDF based solely on Z\_ecdh. This ensures that a PQC-capable SM-DP+ can still provision legacy eUICCs, and a legacy SM-DP+ can still provision PQC-capable eUICCs.

\subsection{Protocol Extensions}

To transport ML-KEM key material within SGP.22 messages, we define two custom ASN.1 TLV tags as shown in Table \ref{tab:asn1-extensions}. These tags are allocated from the context-specific class to avoid conflicts with existing SGP.22 tags.

\begin{table}[h]
\centering
\caption{ASN.1 TLV Extensions for PQC Support}
\label{tab:asn1-extensions}
\begin{tabular}{@{}llp{6cm}@{}}
\toprule
\textbf{Tag} & \textbf{Hex} & \textbf{Description} \\
\midrule
euiccPkKem & 0x5F4A & ML-KEM-768 public key (1,184 bytes) transmitted in PrepareDownloadResponse within euiccSigned2 SEQUENCE. Optional field, present only if eUICC supports PQC. \\
\midrule
smdpCtKem & 0x5F4B & ML-KEM-768 ciphertext (1,088 bytes) transmitted in InitialiseSecureChannelRequest. Optional field, present only if SM-DP+ detects eUICC PQC capability and enters hybrid mode. \\
\bottomrule
\end{tabular}
\end{table}

The PrepareDownloadResponse structure is extended as follows in ASN.1 notation:

\begin{verbatim}
PrepareDownloadResponse ::= [33] SEQUENCE {
  downloadResponseOk [0] SEQUENCE {
    euiccSigned2 [APPLICATION 30] SEQUENCE {
      transactionId [0] OCTET STRING,
      euiccOtpk [APPLICATION 73] OCTET STRING,
      smdpOid [APPLICATION 75] OBJECT IDENTIFIER,
      euiccPkKem [APPLICATION 74] OCTET STRING OPTIONAL
    },
    euiccSignature2 [APPLICATION 55] OCTET STRING
  }
}
\end{verbatim}

The InitialiseSecureChannelRequest structure is extended similarly:

\begin{verbatim}
InitialiseSecureChannelRequest ::= [35] SEQUENCE {
  remoteOpId [0] INTEGER,
  transactionId [0] OCTET STRING,
  controlRefTemplate [6] SEQUENCE {
    keyType [0] OCTET STRING,
    keyLen [1] INTEGER
  },
  smdpOtpk [APPLICATION 73] OCTET STRING,
  smdpCtKem [APPLICATION 75] OCTET STRING OPTIONAL,
  smdpSign [APPLICATION 55] OCTET STRING
}
\end{verbatim}

The use of long-form length encoding is required for the ML-KEM fields. In BER/DER encoding, if the content length exceeds 127 bytes, the length field uses multiple bytes: the first byte has the high bit set and indicates the number of subsequent length bytes, followed by the actual length in big-endian format. For a 1,184-byte ML-KEM public key, the encoding is: tag (0x5F 0x4A), length (0x82 0x04 0xA0, indicating 1,184 in two bytes), followed by 1,184 bytes of key data. Implementations must correctly parse this long-form encoding and allocate sufficient buffer space.

\section{Implementation Details}

\subsection{eUICC-side Implementation}

We implemented the eUICC functionality as a modification to the v-euicc virtual eUICC codebase, written in C. The v-euicc provides a software simulation of an eUICC that communicates via ISO 7816-4 APDU commands over a TCP socket, enabling testing without physical hardware. Our modifications integrate the liboqs library (version 0.11.0) to provide ML-KEM-768 operations.

The core modification is in the PrepareDownload command handler. When the handler receives a PrepareDownloadRequest APDU (tag 0xBF21), it checks the eUICC's PQC capability flags stored in the global state structure. If \texttt{mlkem768\_supported} is true, the handler invokes the hybrid keypair generation routine shown in Algorithm \ref{alg:keygen}.

\begin{algorithm}[h]
\caption{Hybrid Keypair Generation (eUICC)}
\label{alg:keygen}
\begin{algorithmic}[1]
\REQUIRE PQC capability flag \texttt{mlkem768\_supported} = true
\ENSURE Classical ECDH keypair (otPK, otSK) and ML-KEM keypair (pk\_kem, sk\_kem)
\STATE Generate ECDH keypair using NIST P-256 curve
\STATE otPK $\leftarrow$ Extract public key (65 bytes, uncompressed point)
\STATE otSK $\leftarrow$ Extract private key (32 bytes, scalar)
\IF{PQC enabled}
    \STATE kem $\leftarrow$ OQS\_KEM\_new("ML-KEM-768")
    \STATE ret $\leftarrow$ OQS\_KEM\_keypair(kem, pk\_kem, sk\_kem)
    \IF{ret $\neq$ OQS\_SUCCESS}
        \STATE Log error and fall back to classical mode
        \STATE Set \texttt{hybrid\_mode\_active} $\leftarrow$ false
    \ELSE
        \STATE Set \texttt{hybrid\_mode\_active} $\leftarrow$ true
        \STATE Store pk\_kem (1,184 bytes) and sk\_kem (2,400 bytes) in state
    \ENDIF
    \STATE OQS\_KEM\_free(kem)
\ENDIF
\RETURN otPK, otSK, pk\_kem (if hybrid mode), sk\_kem (if hybrid mode)
\end{algorithmic}
\end{algorithm}

The generated ML-KEM public key is added to the euiccSigned2 structure as a TLV field with tag 0x5F4A. A critical implementation challenge was buffer sizing. The original v-euicc code allocated a 256-byte static buffer for euiccSigned2 content, which overflowed when adding the 1,184-byte ML-KEM public key plus TLV overhead (total ~1,190 bytes). We increased this buffer to 2,048 bytes to accommodate the additional data. Failure to do so resulted in silent memory corruption and incorrect signature generation.

The InitialiseSecureChannel command handler (tag 0xBF23) was modified to parse the optional ML-KEM ciphertext field. If present, the handler invokes ML-KEM-768 decapsulation to recover the shared secret Z\_mlkem. The decapsulation call is:

\begin{verbatim}
ret = OQS_KEM_decaps(kem, shared_secret_mlkem, 
                     ciphertext, sk_kem);
\end{verbatim}

After successfully obtaining both Z\_ecdh (from ECDH) and Z\_mlkem (from ML-KEM decapsulation), the handler calls the hybrid KDF function to derive the session keys. The ML-KEM secret key is securely wiped from memory immediately after decapsulation using \texttt{memset(sk\_kem, 0, sk\_kem\_len)} and then freed, as it is no longer needed for subsequent operations. This minimizes the window during which the secret key resides in memory and reduces exposure to side-channel attacks.

\subsection{SM-DP+-side Implementation}

The SM-DP+ functionality is implemented in Python as part of the osmo-smdpp simulator. We extended this codebase with a \texttt{hybrid\_ka.py} module that encapsulates PQC operations using the liboqs-python bindings (version 0.11.0). When the SM-DP+ receives a PrepareDownloadResponse from the eUICC, it parses the euiccSigned2 structure and searches for the optional euiccPkKem field (tag 0x5F4A). If found, the SM-DP+ instantiates a HybridKeyAgreement object and enters hybrid mode.

During BPP generation, the SM-DP+ performs ECDH to obtain Z\_ecdh, then invokes ML-KEM encapsulation:

\begin{verbatim}
ct_kem, ss_mlkem = hybrid_ka.encapsulate(euicc_pk_kem)
\end{verbatim}

The encapsulate method internally calls \texttt{oqs.KeyEncapsulation("ML-KEM-768")} and returns the ciphertext (1,088 bytes) and shared secret (32 bytes). The ciphertext is inserted into the InitialiseSecureChannelRequest message as a TLV field with tag 0x5F4B. Both the SM-DP+ and eUICC then independently invoke the hybrid KDF using their respective copies of Z\_ecdh and Z\_mlkem.

\subsection{Hybrid Key Derivation Function}

The hybrid KDF must combine two independent shared secrets in a manner that preserves security even if one is compromised. We employ a nested HKDF-Extract construction as described in Algorithm \ref{alg:hybrid-kdf}. This approach applies domain separation by labeling each shared secret with its algorithm name before extraction, then combines the resulting pseudo-random keys through hashing.

\begin{algorithm}[h]
\caption{Hybrid Session Key Derivation}
\label{alg:hybrid-kdf}
\begin{algorithmic}[1]
\REQUIRE Z\_ecdh (32 bytes), Z\_mlkem (32 bytes)
\ENSURE KEK (16 bytes), KM (16 bytes)
\STATE label\_ecdh $\leftarrow$ "ECDH-P256"
\STATE label\_mlkem $\leftarrow$ "ML-KEM-768"
\STATE PRK\_ecdh $\leftarrow$ HKDF-Extract(salt=0, IKM=Z\_ecdh || label\_ecdh)
\STATE PRK\_mlkem $\leftarrow$ HKDF-Extract(salt=0, IKM=Z\_mlkem || label\_mlkem)
\STATE Combined\_PRK $\leftarrow$ SHA-256(PRK\_ecdh || PRK\_mlkem)
\FOR{counter $i$ from 1 to 2}
    \STATE T\_prev $\leftarrow$ empty (if $i$ = 1), else previous T
    \STATE T $\leftarrow$ SHA-256(T\_prev || Combined\_PRK || counter\_byte($i$))
    \IF{$i$ = 1}
        \STATE KEK $\leftarrow$ T[0:16]
    \ELSE
        \STATE KM $\leftarrow$ T[0:16]
    \ENDIF
\ENDFOR
\RETURN KEK, KM
\end{algorithmic}
\end{algorithm}

The HKDF-Extract step uses HMAC-SHA-256 with a zero salt and the concatenation of the shared secret with its domain label as input key material. This produces a 32-byte pseudo-random key for each component. These PRKs are concatenated and hashed to form a combined PRK. The final KEK and KM are derived using a simple counter-based KDF: hash the combined PRK with counter 1 to get KEK (first 16 bytes), hash the combined PRK with counter 2 to get KM (first 16 bytes). This construction ensures that compromise of either Z\_ecdh or Z\_mlkem alone does not compromise the derived keys, as an attacker would need to find preimages under SHA-256 to recover the combined PRK.

\subsection{APDU Segmentation for Large Payloads}

The ISO 7816-4 APDU protocol used for eUICC communication typically limits command and response sizes. Standard APDUs support up to 255 bytes of data in short-form encoding and up to 65,535 bytes in extended-length encoding. However, many LPA implementations and eUICC drivers do not fully support extended-length APDUs and instead rely on APDU chaining to transmit large payloads.

In APDU chaining, a large message is split into multiple command APDUs, each marked with a chaining flag in the class byte (CLA). The eUICC buffers these segments and processes the complete command only after receiving the final APDU (with chaining flag cleared). Our implementation in v-euicc supports APDU chaining with a configurable segment buffer size. We increased the default buffer capacity from 512 bytes to 4,096 bytes to accommodate the InitialiseSecureChannelRequest, which can exceed 1,300 bytes when including the ML-KEM ciphertext.

The segmentation is handled transparently by the LPA. When transmitting a PrepareDownloadResponse containing a 1,355-byte payload, the LPA automatically splits it into multiple APDUs of 120 bytes each (typical for socket-based transport), totaling 12 segments. The eUICC reassembles these segments before parsing the command. This mechanism works correctly for both classical and hybrid modes without modification, as it operates at the transport layer below the protocol message layer.

\section{Experimental Setup}

\subsection{Testbed Architecture}

We constructed a virtual testbed to evaluate the performance and correctness of the hybrid key agreement protocol. The testbed comprises three software components:

\textbf{v-euicc}: A C-based virtual eUICC implementation that simulates eUICC functionality including certificate storage, APDU command processing, cryptographic operations, and profile installation. We modified v-euicc to integrate liboqs 0.11.0 for ML-KEM operations and added instrumentation for performance profiling. The v-euicc daemon listens on TCP port 8765 and accepts APDU commands.

\textbf{osmo-smdpp}: A Python-based SM-DP+ simulator that implements the server-side logic of SGP.22. We extended osmo-smdpp with the hybrid\_ka.py module using liboqs-python 0.11.0 to perform ML-KEM encapsulation. The simulator serves profile packages over HTTP on port 8000, proxied through nginx on port 8443 with TLS termination.

\textbf{lpac}: An open-source LPA client implementation in C that communicates with the eUICC via APDU drivers and with the SM-DP+ via HTTPS. We used lpac unmodified as it operates transparently at the transport layer and does not parse the cryptographic contents of PrepareDownload or InitialiseSecureChannel messages.

The testbed configuration is summarized in Table \ref{tab:testbed-config}. All components ran on a single development machine (Apple M1 Pro, 16GB RAM, macOS 15.2) to eliminate network latency variability and enable precise performance measurement.

\begin{table}[h]
\centering
\caption{Testbed Software Configuration}
\label{tab:testbed-config}
\begin{tabular}{@{}llp{5.5cm}@{}}
\toprule
\textbf{Component} & \textbf{Version} & \textbf{Build Configuration} \\
\midrule
v-euicc & Git master (2024-11-10) & CMake with -DENABLE\_PQC=ON, linked against liboqs 0.11.0, OpenSSL 3.4.0 \\
\midrule
osmo-smdpp & Git master (2024-11-10) & Python 3.12.7, liboqs-python 0.11.0 \\
\midrule
lpac & v1.2.0 & CMake default configuration, libcurl 8.11.0 \\
\midrule
liboqs & 0.11.0 & Homebrew installation, ML-KEM-768 implementation from NIST Round 4 submission \\
\midrule
OpenSSL & 3.4.0 & Homebrew installation, used for ECDH and ECDSA operations \\
\bottomrule
\end{tabular}
\end{table}

\subsection{Test Scenarios}

We evaluated two scenarios:

\textbf{Classical Mode (Baseline)}: The eUICC does not advertise PQC support. PrepareDownloadResponse contains only the classical ECDH public key. InitialiseSecureChannelRequest contains only the SM-DP+ ECDH public key. Session keys are derived using the standard SGP.22 KDF based solely on the ECDH shared secret. This represents the current state-of-the-art for RSP provisioning.

\textbf{PQC Hybrid Mode}: The eUICC advertises ML-KEM-768 support in EUICCInfo2. PrepareDownloadResponse contains both ECDH and ML-KEM public keys. The SM-DP+ detects the ML-KEM key and enters hybrid mode. InitialiseSecureChannelRequest contains ECDH public key and ML-KEM ciphertext. Session keys are derived using the hybrid KDF combining both ECDH and ML-KEM shared secrets. This provides quantum resistance while maintaining classical security as a fallback.

For each scenario, we executed a complete profile download and installation flow for the test profile "TS48V2-SAIP2-1-BERTLV-UNIQUE" (a GSMA test profile with known characteristics). The flow includes all four protocol phases: authentication, PrepareDownload, BPP download with 18 APDU commands, and installation result. We repeated each scenario 5 times and report mean values.

\subsection{Metrics Collection Methodology}

We instrumented the v-euicc codebase with high-resolution timing measurements using \texttt{clock\_gettime(CLOCK\_MONOTONIC)} to capture microsecond-precision timestamps before and after each cryptographic operation. The profiling code logs the operation name and elapsed time to stderr:

\begin{verbatim}
[PROFILE] mlkem_keypair: 662 µs (0.662 ms)
[PROFILE] mlkem_decaps: 27 µs (0.027 ms)
[PROFILE] hybrid_kdf: 83 µs (0.083 ms)
\end{verbatim}

Message sizes were extracted by logging the final byte count of each protocol message after TLV encoding:

\begin{verbatim}
[METRICS] PrepareDownloadResponse: 1355 bytes
\end{verbatim}

We developed a post-processing script (\texttt{extract-true-metrics.sh}) that parses the eUICC and SM-DP+ log files, extracts timing and size measurements using regular expressions, and outputs structured metrics in key-value format. These metrics were then fed into a Python script (\texttt{generate-visualizations.py}) that produced CSV exports and comparative charts using matplotlib.

Critically, all measurements are derived from actual execution of the testbed. We do not simulate timing or fabricate data. The profiling instrumentation introduces negligible overhead (< 1 microsecond per measurement) compared to the cryptographic operations being measured, ensuring that the reported values accurately reflect real performance characteristics.

\section{Evaluation Results}

\subsection{Key Size Overhead}

Figure \ref{fig:key-sizes} compares the key material sizes for classical and hybrid modes. The classical mode uses only ECDH with a 65-byte public key and 32-byte private key. The hybrid mode adds ML-KEM-768 with a 1,184-byte public key, 2,400-byte secret key, and 1,088-byte ciphertext.

\begin{figure}[h]
  \centering
  \includegraphics[width=0.8\columnwidth]{output/chart_key_sizes.png}
  \caption{Comparison of cryptographic key sizes between classical ECDH-only and PQC hybrid modes. ML-KEM-768 introduces 2,272 bytes of additional key material transmitted over the air (public key + ciphertext).}
  \label{fig:key-sizes}
\end{figure}

Table \ref{tab:key-sizes} provides detailed breakdowns. The total PQC overhead is 2,272 bytes (public key + ciphertext), representing the additional data that must be transmitted during provisioning. This overhead is incurred exactly once per profile installation. Given that a typical eSIM profile is 50-200 KB in size (including the encrypted BPP), the 2.3 KB overhead represents a 1-5\% increase in total provisioning data. For a one-time provisioning operation, this is acceptable.

\begin{table}[h]
\centering
\caption{Detailed Key Size Comparison}
\label{tab:key-sizes}
\begin{tabular}{@{}lrr@{}}
\toprule
\textbf{Key Material} & \textbf{Classical (bytes)} & \textbf{PQC Hybrid (bytes)} \\
\midrule
ECDH Public Key & 65 & 65 \\
ECDH Private Key & 32 & 32 \\
ML-KEM Public Key & 0 & 1,184 \\
ML-KEM Secret Key & 0 & 2,400 \\
ML-KEM Ciphertext & 0 & 1,088 \\
\midrule
\textbf{Total Transmitted} & \textbf{65} & \textbf{2,337} \\
\textbf{Total Stored (eUICC)} & \textbf{97} & \textbf{2,497} \\
\bottomrule
\end{tabular}
\end{table}

The 2,400-byte ML-KEM secret key poses a memory challenge for resource-constrained eUICC implementations. A typical JavaCard-based eUICC has 256 KB to 1 MB of non-volatile memory and 8 KB to 64 KB of transient RAM. Storing the ML-KEM secret key requires 2.4 KB, which is manageable but non-trivial. Our implementation mitigates this by using transient storage: the secret key is allocated during PrepareDownload, used for a single decapsulation in InitialiseSecureChannel (typically within 1-2 seconds), and then immediately wiped and freed. This minimizes memory footprint and exposure to attacks.

\subsection{Protocol Message Size Analysis}

Figure \ref{fig:message-sizes} and Table \ref{tab:message-sizes} show the impact of PQC on individual protocol message sizes. The PrepareDownloadResponse increases from 162 bytes to 1,355 bytes (8.4x), driven almost entirely by the addition of the 1,184-byte ML-KEM public key. Similarly, InitialiseSecureChannelRequest increases from 178 bytes to 1,272 bytes (7.1x) due to the 1,088-byte ML-KEM ciphertext.

\begin{figure}[h]
  \centering
  \includegraphics[width=0.8\columnwidth]{output/chart_message_sizes.png}
  \caption{Protocol message size comparison. PrepareDownload and InitialiseSecureChannel messages exhibit 7-8x size increases due to ML-KEM key material. BoundProfilePackage size is dominated by profile data and shows minimal increase.}
  \label{fig:message-sizes}
\end{figure}

\begin{table}[h]
\centering
\caption{Protocol Message Size Breakdown}
\label{tab:message-sizes}
\begin{tabular}{@{}lrrr@{}}
\toprule
\textbf{Message} & \textbf{Classical} & \textbf{PQC} & \textbf{Increase} \\
\midrule
PrepareDownloadResponse & 162 B & 1,355 B & 8.4x \\
BoundProfilePackage (BF36 wrapper) & 183 B & 1,277 B & 7.0x \\
InitialiseSecureChannelRequest & 178 B & 1,272 B & 7.1x \\
Total Profile Data (encrypted) & 1,689 B & 1,689 B & 1.0x \\
\bottomrule
\end{tabular}
\end{table}

Notably, the total encrypted profile data remains unchanged at 1,689 bytes because the profile content itself is not affected by the key exchange mechanism. The BPP encryption uses AES-128-CBC with keys derived from the session key establishment. Whether those session keys are derived classically or via hybrid mode does not impact the size of the encrypted data. This demonstrates that PQC overhead is isolated to the key establishment phase and does not amplify throughout the protocol.

\subsection{Performance Overhead}

Figure \ref{fig:performance} shows the computational overhead of PQC operations. ML-KEM keypair generation takes 0.128 ms, decapsulation takes 0.026 ms, and the hybrid KDF takes 0.078 ms, for a total PQC overhead of 0.232 ms per provisioning session. These measurements were obtained on an Apple M1 Pro processor (3.2 GHz) and represent single-threaded performance.

\begin{figure}[h]
  \centering
  \includegraphics[width=0.8\columnwidth]{output/chart_performance.png}
  \caption{Performance timing breakdown for PQC operations. Total computational overhead is 0.232 ms, negligible compared to typical network latency (50-200 ms) and total provisioning time (5-10 seconds).}
  \label{fig:performance}
\end{figure}

To contextualize these numbers, consider that a typical profile download session over a 4G/5G cellular connection involves:
\begin{itemize}
\item Network round-trip time: 50-200 ms
\item TLS handshake: 100-300 ms
\item BPP download time: 500-2000 ms (for 50-200 KB profile at 100-200 Kbps effective throughput)
\item Profile installation time: 1000-3000 ms (cryptographic operations, file system writes)
\end{itemize}

The 0.232 ms PQC overhead is less than 0.5\% of the total provisioning time and is completely dominated by network latency and profile installation overhead. Even on lower-end embedded processors (e.g., ARM Cortex-M4 at 100 MHz), ML-KEM operations scale linearly and would complete in under 10 ms, still negligible relative to network operations.

\subsection{Bandwidth Impact}

Figure \ref{fig:bandwidth} illustrates the total bandwidth consumption across the entire provisioning flow. Classical mode consumes 523 bytes total (162 bytes upload in PrepareDownload + 361 bytes download in BPP), while PQC hybrid mode consumes 3,904 bytes (1,355 bytes upload + 2,549 bytes download), representing a 7.5x increase (646\% overhead).

\begin{figure}[h]
  \centering
  \includegraphics[width=0.8\columnwidth]{output/chart_bandwidth.png}
  \caption{Total bandwidth comparison. PQC introduces 3,381 additional bytes (646\% increase), concentrated in PrepareDownload and InitialiseSecureChannel messages. This represents a one-time cost per profile installation.}
  \label{fig:bandwidth}
\end{figure}

At first glance, a 7.5x bandwidth increase may appear significant. However, this must be evaluated in the context of real-world provisioning:

\begin{itemize}
\item The additional 3.4 KB is transmitted once per profile installation. eSIM profiles are typically installed infrequently (once per device activation, or occasionally when changing carriers).
\item The 3.4 KB represents 2-7\% overhead relative to the total profile download size (50-200 KB). The vast majority of bandwidth is consumed by the encrypted profile data, which is unaffected by PQC.
\item At typical cellular data rates (1-10 Mbps), transmitting an extra 3.4 KB adds 3-30 milliseconds of transfer time, negligible compared to network latency variability.
\item For IoT devices on constrained networks (e.g., NB-IoT at 50 Kbps), the 3.4 KB adds ~500 ms. This is measurable but acceptable for provisioning operations that occur rarely.
\end{itemize}

In summary, while the relative bandwidth overhead is substantial (7.5x), the absolute overhead is small enough that it does not pose a practical barrier to deployment. The one-time cost of quantum resistance is justified by the long-term security benefits, especially for devices with multi-year operational lifetimes.

\section{Discussion: Deployment Challenges and Insights}

\subsection{Implementation Challenges Encountered}

During implementation and testing, we encountered several practical challenges that highlight the complexity of integrating PQC into legacy protocols:

\textbf{Buffer Overflow Vulnerability.} The most critical issue was a buffer overflow in the PrepareDownload response construction. The original v-euicc code allocated a 256-byte stack buffer for the euiccSigned2 content, assuming that all fields (transactionId, euiccOtpk, smdpOid) would fit comfortably. When we added the 1,184-byte ML-KEM public key, the code continued writing beyond the buffer boundary, corrupting adjacent stack memory. This resulted in unpredictable behavior including crashes and incorrect ECDSA signature generation (which failed verification at the SM-DP+). The fix required increasing the buffer to 2,048 bytes and adding explicit bounds checking before memory copies. Similarly, the prep\_resp\_buf for the full PrepareDownloadResponse was increased from 512 to 2,048 bytes. This experience underscores the importance of defensive programming and robust buffer management when integrating large PQC keys into protocols designed for classical cryptography.

\textbf{TLV Long-Form Length Encoding.} SGP.22 uses ASN.1 BER/DER encoding for all protocol messages. In DER, lengths up to 127 bytes use short-form encoding (single byte), while larger lengths require long-form encoding. For the 1,184-byte ML-KEM public key, the proper encoding is: tag 0x5F4A (2 bytes), length 0x82 0x04 0xA0 (3 bytes: 0x82 indicates 2-byte length follows, 0x04A0 = 1,184), then 1,184 bytes of data. We initially overlooked the long-form encoding, causing TLV parsing failures at the SM-DP+. The fix required updating the build\_tlv() function to automatically select short-form or long-form encoding based on the content length and updating all TLV parsers to handle both forms correctly.

\textbf{NULL Pointer Dereference in Decapsulation Path.} In early testing, the v-euicc crashed with SIGSEGV during ML-KEM decapsulation. Debugging revealed that the InitialiseSecureChannelRequest parser successfully extracted the ciphertext pointer but failed to validate it before passing to mlkem\_decapsulate(). If the TLV parsing encountered malformed data, it could return a NULL pointer, which was then dereferenced inside liboqs. We added explicit NULL checks and length validation before invoking any liboqs functions: \texttt{if (ciphertext \&\& ct\_len == 1088) ...}. This is a general lesson: when integrating external cryptographic libraries, always validate inputs at trust boundaries.

\textbf{Dynamic Library Loading on macOS.} During Python integration of liboqs-python, we encountered issues with the dynamic linker failing to locate liboqs.dylib. macOS System Integrity Protection (SIP) strips DYLD\_LIBRARY\_PATH from environment variables for security, breaking the usual library loading mechanism. We resolved this by: (1) installing liboqs to a standard system location (/opt/homebrew/lib), (2) monkey-patching ctypes.util.find\_library() to explicitly search common paths, and (3) ensuring the Python virtual environment properly propagates library paths. This is a deployment challenge specific to certain operating systems but highlights the need for careful dependency management.

\subsection{Backward Compatibility Strategy}

A critical design goal was ensuring that PQC-capable eUICCs can interoperate with legacy SM-DP+ servers, and vice versa. We achieve this through capability-based negotiation:

\begin{itemize}
\item The eUICC advertises its PQC support by including a custom capability flag in the EUICCInfo2 structure returned during chip info requests. If the eUICC supports ML-KEM-768, it sets this flag.
\item The eUICC includes the ML-KEM public key in PrepareDownloadResponse only if it supports PQC. Legacy SM-DP+ servers that do not recognize tag 0x5F4A simply ignore the unknown field (per ASN.1 forward compatibility rules) and proceed with classical ECDH.
\item The SM-DP+ checks for the presence of tag 0x5F4A in the PrepareDownloadResponse. If found, it enters hybrid mode. If absent, it uses classical mode. This decision is made per-session, allowing the same SM-DP+ to serve both legacy and PQC-capable eUICCs.
\item If a legacy eUICC receives an InitialiseSecureChannelRequest containing tag 0x5F4B (ML-KEM ciphertext), it ignores the unknown field and uses only the ECDH public key for key derivation.
\end{itemize}

This bilateral capability negotiation ensures that all four combinations (legacy eUICC + legacy SM-DP+, legacy eUICC + PQC SM-DP+, PQC eUICC + legacy SM-DP+, PQC eUICC + PQC SM-DP+) interoperate correctly. The first three cases fall back to classical mode, while only the fourth case uses hybrid mode. This enables gradual deployment across the ecosystem without requiring synchronized upgrades.

We validated backward compatibility by compiling two versions of v-euicc: one with -DENABLE\_PQC=OFF (classical-only) and one with -DENABLE\_PQC=ON (hybrid-capable). We then ran the classical version against the PQC-enabled SM-DP+ and confirmed successful profile installation using purely classical cryptography, with no errors or protocol failures. This demonstrates that the protocol extensions are truly optional and do not break existing implementations.

\subsection{Memory Constraints and Secure Key Management}

The 2,400-byte ML-KEM secret key poses challenges for memory-constrained eUICC implementations. High-end eUICCs based on JavaCard 3.0.4 or later typically have 512 KB to 1 MB of persistent memory (EEPROM or flash) and 32 KB to 64 KB of transient memory (RAM). Embedded devices and IoT eUICCs may have significantly less, with some constrained designs having only 128 KB persistent storage and 8 KB RAM.

Our implementation addresses this through temporal key management: the ML-KEM keypair is generated in RAM during PrepareDownload processing and stored in a transient buffer. The secret key resides in memory for at most a few seconds (the time between PrepareDownload and InitialiseSecureChannel, typically 1-3 seconds over a network). After decapsulation, the secret key is immediately wiped using explicit\_bzero() or memset\_s() (compiler-barrier-protected memory zeroing functions) and the memory is freed. This minimizes the attack surface for side-channel attacks that attempt to extract the secret key from RAM.

For devices with extremely limited RAM, an alternative approach would be to generate the ML-KEM keypair outside the eUICC (e.g., in the LPA or device operating system) and pass only the public key to the eUICC for inclusion in PrepareDownloadResponse. The secret key would be retained by the LPA, and decapsulation would occur on the LPA side. The eUICC would then receive only the derived session keys. This "offloading" approach trades increased protocol complexity for reduced memory footprint on the eUICC.

\subsection{Standardization Considerations}

For production deployment, several standardization activities are necessary:

\textbf{Formal TLV Tag Allocation.} Our prototype uses custom tags 0x5F4A and 0x5F4B from the context-specific class. For standardization, GSMA should formally allocate these tags (or alternative tags) in the SGP.22 specification and define their structure, encoding, and usage constraints. This prevents conflicts with future protocol extensions.

\textbf{Capability Negotiation Formalization.} The EUICCInfo2 structure should be extended with a formal \texttt{pqcCapabilities} field indicating which PQC algorithms the eUICC supports (e.g., ML-KEM-768, ML-KEM-1024, future algorithms). This allows the SM-DP+ to select an appropriate algorithm based on security policy and device capabilities.

\textbf{Hybrid KDF Specification.} The exact hybrid KDF construction (algorithm, labels, counter format) must be standardized to ensure interoperability between vendors. Different KDF constructions may produce different session keys even from the same input shared secrets, breaking compatibility.

\textbf{Security Level Negotiation.} Future versions may support multiple ML-KEM variants (ML-KEM-512, ML-KEM-768, ML-KEM-1024) corresponding to NIST security levels 1, 3, and 5. The protocol should include negotiation to select the highest mutually supported level.

\textbf{Migration Timeline Coordination.} The NIST PQC standards were finalized in August 2024. GSMA should begin incorporating PQC into SGP.22 v3.0 or later, targeting publication in 2025-2026. This gives eUICC manufacturers and SM-DP+ operators a 2-3 year window to implement PQC before the quantum threat becomes acute (projected 2030-2035). Early adoption is recommended given the long deployment cycles in telecommunications.

\section{Related Work}

Post-Quantum Cryptography integration into real-world protocols has been explored across multiple domains. The TLS 1.3 hybrid key exchange proposals by Stebila et al. combine classical ECDHE with PQC KEMs to provide transitional security during the PQC migration period. These proposals influenced our hybrid KDF design, particularly the use of domain-separated labels to prevent cross-protocol attacks. However, TLS operates at the transport layer with different constraints (higher bandwidth, more computation available, ephemeral connections), whereas RSP operates at the application layer with strict memory limits and long-lived trust relationships.

The ETSI Technical Committee Cyber has published technical reports on quantum-safe cryptography in telecommunications, including recommendations for transitioning mobile networks to PQC. Their work focuses primarily on 5G core network security (authentication and key agreement protocols) rather than eSIM provisioning. Our work complements these efforts by addressing the specific challenges of SGP.22 integration.

In the context of smart cards and secure elements, several works have explored PQC feasibility on constrained devices. Howe et al. demonstrated ML-KEM implementation on ARM Cortex-M4 microcontrollers, achieving keypair generation in under 1 ms. Oder et al. analyzed memory requirements for lattice-based cryptography on JavaCard platforms, concluding that ML-KEM-768 is feasible on high-end cards but requires careful optimization. Our experimental results align with these findings: PQC operations complete quickly (< 1 ms) but impose non-trivial memory overhead (2.4 KB for the secret key).

To our knowledge, this is the first work to implement and evaluate PQC integration specifically for GSMA SGP.22 Remote SIM Provisioning. Previous work has proposed PQC for general IoT provisioning protocols or smart card applications but has not addressed the unique requirements of RSP: the GSMA PKI trust model, ASN.1 BER/DER encoding, APDU-based communication, and backward compatibility constraints.

\section{Conclusion}

This paper presented a practical strategy for integrating Post-Quantum Cryptography into GSMA Consumer Remote SIM Provisioning (SGP.22) through hybrid key exchange combining classical ECDH with ML-KEM-768. Our implementation demonstrates that quantum resistance can be achieved with acceptable overhead: 2,272 bytes of additional key material, 7.5x bandwidth increase, and 0.232 ms of computational overhead per provisioning session. These costs are dominated by network latency and profile installation time in real-world deployments, making PQC integration feasible without significant user experience impact.

The key contributions of this work are: (1) a hybrid key agreement protocol that preserves the GSMA PKI hierarchy and provides backward compatibility with legacy eUICCs and SM-DP+ servers, (2) ASN.1 protocol extensions for transporting ML-KEM key material within existing SGP.22 message structures, (3) an open-source implementation and experimental evaluation on a virtual testbed quantifying real performance characteristics, and (4) documentation of practical deployment challenges including buffer management, TLV encoding, and memory constraints.

Future work will address signature migration from ECDSA to ML-DSA (formerly Dilithium), enabling end-to-end quantum resistance for both confidentiality and authenticity. We also plan to explore optimizations for extremely memory-constrained eUICCs, including offloading PQC operations to the LPA and investigating lightweight hybrid constructions. Formal security analysis of the hybrid KDF construction and integration with automated theorem provers (e.g., ProVerif, Tamarin) would provide additional confidence in the protocol's security properties. Finally, collaboration with GSMA and eUICC manufacturers to standardize the protocol extensions and coordinate deployment timelines remains essential for ecosystem-wide adoption.

The store-now-decrypt-later threat is real and immediate. By demonstrating the feasibility of PQC integration into RSP with minimal disruption and acceptable overhead, this work provides a concrete path toward quantum-resistant eSIM provisioning, protecting subscriber credentials and operator infrastructure against future quantum adversaries.

\bibliographystyle{plain}
\begin{thebibliography}{9}

\bibitem{GSMA_SGP22}
GSM Association (GSMA). 
\emph{SGP.22: Remote SIM Provisioning Architecture for Consumer Devices}. 
GSMA, 2023.

\bibitem{GSMA_SGP24}
GSM Association (GSMA). 
\emph{SGP.24: RSP Technical Specification for Consumer Devices}. 
GSMA, 2024.

\bibitem{NIST_PQC2024}
National Institute of Standards and Technology (NIST). 
\emph{Post-Quantum Cryptography Standardization Project}, 2024. 
Available at: \url{https://csrc.nist.gov/projects/post-quantum-cryptography}

\bibitem{OQS2024}
Open Quantum Safe Project. 
\emph{liboqs and OQS-OpenSSL Libraries}. 
Accessed 2025. 
Available at: \url{https://openquantumsafe.org/}

\bibitem{Kyber2024}
Bos, J. W., Ducas, L., Kiltz, E., Lepoint, T., Lyubashevsky, V., Schwabe, P., Seiler, G., and Stehlé, D. 
\emph{CRYSTALS-Kyber: Algorithm Specifications and Supporting Documentation}. 
NIST PQC Round 4 Submission, 2024.

\end{thebibliography}

\end{document}
